\subsection{A moment condition on maximizing tests}

For a general Hermitian background, a uniform bound on the moments of normalized potentials with respect to the probabilities of one smooth maximizing sequence can replace the regularity of the maximizing current. Together with the envelope condition in a smaller background, it shows that the cluster limits of the solutions of the exponential equations are exact threshold solutions.

For a nonzero $0\leq f\in L^1(\mathrm{d}V)$ with $b_\omega(f)>0$, let $\mathcal S_f$ be the class of threshold subsolutions with supremum zero and $\Phi_f$ its greatest member, from \cref{prop:greatest-normalized-threshold}. As in \cref{prop:semipositive-dual-exponential-limit}, $u_\lambda$ denotes the greatest solution of the exponential equation $\mathrm{MA}^{\mathrm{ext}}_\omega(u)=\mathrm{e}^{\lambda u}f\,\mathrm{d}V$, and $v_\lambda=u_\lambda-\sup_Xu_\lambda$, $c_\lambda=\mathrm{e}^{\lambda\sup_Xu_\lambda}$.

\Needspace{12\baselineskip}
\begin{theorem}\label{thm:gauduchon-moment-threshold}\index{Moment constraint}
Let $0\leq f\in L^1(\mathrm{d}V)$ be nonzero with $b=b_\omega(f)>0$, and put $p=f^{1/n}$ and $R=\rho_\omega(p)=b^{-1/n}$. Suppose that there are smooth strictly positive $T_j\in\mathscr G_\omega$, with associated probability measures $\nu_j$, such that
\[
 d_j=D_p(T_j)\longrightarrow R,\qquad
 \mathrm{d}\nu_j=d_j^{-1}p k_{T_j}\,\mathrm{d}V,
\]
and a constant $C$ satisfying
\begin{equation}\label{eq:maximizing-gauduchon-moments}
 \int_X(-v)\,\mathrm{d}\nu_j\leq C
 \quad\text{for all }j\text{ and all }v\in\PSH(X,\omega)
                  \text{ with }\sup_Xv=0.
\end{equation}
Then, for every $\lambda>0$,
\[
 b\leq c_\lambda\leq b\mathrm{e}^{C\lambda},\qquad
 \Phi_f-C\leq v_\lambda\leq0.
\]
Suppose in addition that $P_{\delta\omega}(\Phi_f)$ is not identically $-\infty$ for some $0<\delta<1$. Whenever $\lambda_k\downarrow0$ and $v_{\lambda_k}\to v$ in $L^1(X)$, the convergence holds in capacity and
\[
 \mathrm{MA}^{\mathrm{ext}}_\omega(v)=bf\,\mathrm{d}V,
 \qquad \Phi_f-C\leq v\leq\Phi_f.
\]
In particular an exact threshold solution exists.
\end{theorem}
\begin{proof}
Since $\mathrm{MA}^{\mathrm{ext}}_\omega(v_\lambda)=c_\lambda\mathrm{e}^{\lambda v_\lambda}f\,\mathrm{d}V$, the trace inequality of \cref{lem:unbounded-gauduchon-trace-density} and Stokes' theorem give, for each $j$ and $\lambda>0$,
\[
 1\geq c_\lambda^{1/n}d_j
              \int_X\mathrm{e}^{\lambda v_\lambda/n}\,\mathrm{d}\nu_j
   \geq c_\lambda^{1/n}d_j\mathrm{e}^{-C\lambda/n}.
\]
\Needspace{9\baselineskip}
The second inequality follows from Jensen's inequality and \eqref{eq:maximizing-gauduchon-moments}. Letting $j\to\infty$ and using $R=b^{-1/n}$, we obtain the upper bound for $c_\lambda$. The lower bound and the comparison with every normalized threshold subsolution were proved in \cref{cor:individual-positive-threshold-exponential}; applying that comparison to $\Phi_f$, we obtain
\[
 v_\lambda\geq\Phi_f-\lambda^{-1}\log(c_\lambda/b)
                \geq\Phi_f-C.
\]

Now let $v_{\lambda_k}\to v$ in $L^1$. By \cref{prop:greatest-normalized-threshold}, the limit $v$ lies in $\mathcal S_f$, hence below $\Phi_f$. The lower bound passes to the limit, so $\Phi_f-C\leq v\leq\Phi_f$; in particular $v\sim\Phi_f$ in the sense of \cref{def:singularity-order}. For all $k$ with $\lambda_k\leq1$,
\[
 \mathrm{MA}^{\mathrm{ext}}_\omega(v_{\lambda_k})
   =c_{\lambda_k}\mathrm{e}^{\lambda_k v_{\lambda_k}}f\,\mathrm{d}V
   \leq b\mathrm{e}^{C}f\,\mathrm{d}V.
\]
The functions $|\mathrm{e}^{v_{\lambda_k}}-\mathrm{e}^v|$ are uniformly bounded and tend to zero in $\mathrm{d}V$-measure, so their integrals against the fixed $L^1$ density $f$ tend to zero. Hence the error term in \eqref{eq:common-barrier-self-error} tends to zero, as required by \cref{prop:lower-barrier-criterion} with $\varphi=\Phi_f$. The additional hypothesis $v\preceq\Phi_f$ of that proposition holds since $v\leq\Phi_f$, and the condition on $P_{\delta\omega}(\Phi_f)$ is assumed. That proposition therefore gives convergence in capacity.

The continuity of damped measures under convergence in capacity, \cref{thm:damped-continuity}, now identifies the equation without a mass-conservation argument. Write $g_k=c_{\lambda_k}\mathrm{e}^{\lambda_k v_{\lambda_k}}f$. We know $g_k\to bf$ in $L^1$ and $0\leq g_k\leq b\mathrm{e}^Cf$ for large $k$. Hence
\[
 \mathrm{e}^{v_{\lambda_k}}\mathrm{MA}^{\mathrm{ext}}_\omega(v_{\lambda_k})
   =\mathrm{e}^{v_{\lambda_k}}g_k\,\mathrm{d}V
   \longrightarrow\mathrm{e}^v bf\,\mathrm{d}V
\]
in total variation, while \cref{thm:damped-continuity} gives the weak limit
$\mathrm{e}^v\mathrm{MA}^{\mathrm{ext}}_\omega(v)$. Dividing by $\mathrm{e}^v$ off the pole set of $v$, which neither measure charges, gives the equation and the finiteness of the mass. Finally, compactness of normalized $\omega$-plurisubharmonic functions provides an $L^1$ cluster limit along any sequence $\lambda_k\downarrow0$, which proves existence.
\end{proof}

The moment condition holds if the densities of the $\nu_j$ are uniformly bounded in $L^q(\mathrm{d}V)$ for some $q>1$, by H\"older's inequality and the uniform bounds for the moments of normalized $\omega$-plurisubharmonic functions. More generally, a uniform entropy bound suffices. Write $a_j=\mathrm{d}\nu_j/\mathrm{d}V$ and suppose $\int_Xa_j\log a_j\,\mathrm{d}V\leq H$ for all $j$. (By \cref{prop:hermitian-dual-entropy-bounds}, a uniform entropy bound is equivalent to a uniform bound for the dual entropy $H_0(\nu_j)$ with zero reference defined there.) Choose $\alpha>0$ and $A$ from \cref{lem:uniform-exponential}, so that $\int_X\mathrm{e}^{-\alpha v}\,\mathrm{d}V\leq A$ for every $v\in\PSH(X,\omega)$ with $\sup_Xv=0$. The elementary entropy inequality, applied first to the truncations $v_T=\max(v,-T)$ at levels $T>0$, gives
\[
 \alpha\int_X(-v_T)a_j\,\mathrm{d}V
 \leq \int_Xa_j\log a_j\,\mathrm{d}V
             +\log\int_X\mathrm{e}^{-\alpha v_T}\,\mathrm{d}V
 \leq H+\log A.
\]
To prove this inequality, apply $t\log t-t+1\geq0$ to $t=a_j/q$, integrate against $q\,\mathrm{d}V$, where $q=\mathrm{e}^{-\alpha v_T}/\int_X\mathrm{e}^{-\alpha v_T}\,\mathrm{d}V$ is a probability density, and rearrange. Monotone convergence as $T\to\infty$ proves \eqref{eq:maximizing-gauduchon-moments} with $C=(H+\log A)/\alpha$.

The same entropy bound also gives attainment of the dual supremum along this sequence. Indeed, $t\log t\geq-\mathrm{e}^{-1}$ and, for $K>1$,
\[
 \int_{\{a_j>K\}}a_j\,\mathrm{d}V
       \leq\frac{H+\mathrm{e}^{-1}}{\log K}.
\]
This proves uniform integrability of the $a_j$, and since $d_j\to R>0$, also of the densities $p k_{T_j}=d_ja_j$. By \cref{prop:selected-gauduchon-attainment}, every weak cluster limit $T$ of $(T_j)$ maximizes $D_p$, and along a subsequence converging to $T$ the probabilities $\nu_j$ converge to $R^{-1}p k_T\,\mathrm{d}V$. For each normalized $v$, the function $-v$ is nonnegative and lower semicontinuous, so the lower semicontinuity of its integral under weak convergence of measures passes the moment bound to this limiting probability. These conclusions give no regularity of $T$.

The envelope condition holds if $\Phi_f$ is bounded. In general, it and the moment bound are additional hypotheses; the tail condition on $f^{1/n}$ alone does not provide them in this argument. The theorem identifies the equation satisfied by each cluster limit. Uniqueness, which gave convergence of the whole family in \cref{prop:semipositive-dual-exponential-limit}, is not available here.
